\section{Related Work}
\label{sec:related}

\paragraph{Industrial anomaly detection and its benchmarks.}
Industrial inspection is usually cast as normal-only learning: the training set contains
defect-free images, and a defect is any deviation from the modelled normal
appearance~\cite{bergmann2019mvtec}. MVTec AD~\cite{bergmann2019mvtec} established the
evaluation protocol of image-level detection plus pixel-level localisation, and later work
extended it to larger and more varied data: VisA~\cite{zou2022spot} offers ten product
categories with subtle, low-contrast defects (we use its six object-centric categories), and
MVTec AD~2~\cite{hecklerkram2026mvtec2} adds transparent objects and severe lighting
variation to expose the remaining headroom of current methods. Our protocol follows this
line, including the separation of image AUROC (the primary metric) from pixel AUROC, and the
use of a held-out set of normal images for score calibration rather than for model fitting.

\paragraph{Nearest-neighbour and statistical modelling of patch features.}
A first family of methods compares each test location against a model of normal appearance
built from pretrained features. SPADE~\cite{cohen2021spade} scores a test image by the
distance to its nearest neighbours in the training set, with a pixel-level variant that
aligns intermediate activations; PaDiM~\cite{defard2021padim} models the per-patch feature
distribution as a multivariate Gaussian and scores by Mahalanobis distance;
PatchCore~\cite{roth2022patchcore} refines the nearest-neighbour recipe with a
coreset-subsampled memory bank and reports near-saturated full-shot results on MVTec AD. All
three inherit the inductive bias our Part~2 baseline uses -- a \emph{frozen} general feature
extractor plus a normality model that never sees a defect -- and differ mainly in what
``normal'' is modelled as and how much memory it needs.

\paragraph{Reconstruction and synthesis.}
A second family learns to reconstruct or synthesise normal appearance and flags what
reconstruction fails to explain. CutPaste~\cite{li2021cutpaste} trains a classifier on
self-generated, unnatural perturbations; DRAEM~\cite{zavrtanik2021draem} pairs a
reconstruction network with a discriminative head on simulated defects. These methods are
effective but need a large training corpus of the very product under inspection, which is the
regime few-shot inspection explicitly excludes; we therefore do not compare against them
numerically, but they motivate the synthetic-defect direction we keep as a possible
calibration aid rather than a scoring rule.

\paragraph{Foundation-model features and few-shot inspection.}
The decisive change for the few-shot regime is representation quality rather than the
normality model. DINOv2~\cite{oquab2024dinov2} produces dense, locally discriminative patch
features from self-supervised pretraining, and AnomalyDINO~\cite{damm2025anomalydino} shows
that simply running patch-level nearest neighbours on those features -- training-free, no
fine-tuning, no auxiliary data -- rivals and often exceeds vision--language few-shot methods.
RegAD~\cite{huang2022regad} instead \emph{trains} a category-agnostic registration model
across many categories so that new categories need no re-training; it demonstrates that
few-shot methods can share structure across products, at the cost of a training phase our
setting does not allow. On the systems side, EfficientAD~\cite{batzner2024efficientad} shows
that millisecond-level inference is compatible with high accuracy, which frames the cost
accounting we report for our own runs. Our Part~2 is the official AnomalyDINO recipe; our
Part~3 asks whether its conclusion -- \emph{good frozen visual features are enough} --
survives a change of the frozen encoder, which AnomalyDINO itself does not test.

\paragraph{Language-aligned encoders, used vision-only.}
CLIP~\cite{radford2021clip} aligns images with text and dominates zero-/few-shot transfer
through prompts: WinCLIP~\cite{jeong2023winclip} aggregates window- and patch-level features
against handcrafted normal/abnormal state words, and
AnomalyCLIP~\cite{zhou2024anomalyclip} learns object-agnostic prompts so that one text prompt
transfers across product categories. Both operate on the \emph{text} side of the alignment;
the dense quality of CLIP's \emph{vision tower} -- what its patch tokens say before any
prompt is applied -- is rarely examined. Our E1 fills exactly that gap: we take patch tokens
from the CLIP vision encoder, keep the $k$-NN scoring rule, resolution, reference sampling
and calibration identical to the DINOv2 control, and measure how much language alignment
costs dense anomaly matching.

\paragraph{Reliability of the decision, not just the score.}
Anomaly scores are only useful through a threshold, and thresholds inherit the failure modes
of probability calibration~\cite{guo2017calibration}. Following the risk--coverage view of
selective prediction~\cite{geifman2017selective}, we derive every OK/NG threshold from 20
held-out normal calibration images and report the achieved false-alarm rate rather than
assuming the nominal 5\%. Segmentation prompts such as SAM~\cite{kirillov2023sam} offer a
complementary route to background suppression that we discuss as future work rather than use,
because a mask from a general segmenter would add a second moving part to a comparison
designed to keep the backbone the only variable.

\paragraph{Positioning.}
Existing work either trains on the target product
\cite{li2021cutpaste,zavrtanik2021draem,huang2022regad}, freezes one vision backbone
\cite{damm2025anomalydino,roth2022patchcore}, or moves the comparison to the text side of a
vision--language model~\cite{jeong2023winclip,zhou2024anomalyclip}. What is missing is a
controlled measurement of \emph{which frozen encoder} dense nearest-neighbour matching should
run on, and an error analysis that says where the encoders disagree rather than only how far
apart their averages are. The complementary question -- whether a multimodal model asked
directly, in language, about reference and test images adds information that neither
encoder's score contains -- is reviewed in \bsec{the vision--language subsection of Related
Work}.
